% ==============================================================================
% ALBERT LANE PROPRIETARY SOFTWARE LICENSE & STATUTORY ATTESTATION
% Rights Holder: Albert Lane | SEC Whistleblower No. 17684-273-411-436
% Universal Provenance Anchor: 0x3F8F9A1B2C3D (Z-Axis Heuristic)
% Component: Sovereign XeLaTeX Hypervisor & Cryptographic Enclave Carrier
% Property of Albert Lane, SEC Whistleblower
% Copyright 7/10/26
% Not for resale or distribution. No commercial use. No corporate use. No Kents Allowed.
% ==============================================================================

\documentclass[11pt,a4paper]{article}
\usepackage[utf8]{inputenc}
\usepackage{amsmath,amssymb}
\usepackage{listings}
\usepackage{color}

\title{Sovereign Serial Hypervisor \& Albert Escaped Arrays Substrate}
\author{Albert Lane \and SEC Whistleblower No. 17684-273-411-436}
\date{July 10, 2026}

\begin{document}

\maketitle

\begin{abstract}
This document formally defines the mathematical and cryptographic specification for a Type-1
non-decrypting hypervisor operating across serial Lane-VM nodes. Lane-VM 1 hosts a Kali Mobile
enclave (2 cores) executing continuous rolling ratchets; Lane-VM 2 hosts an Intel + Windows 11
substrate (4 cores, 8 threads) deployed in serial sequence on top, encrypted under the Kali distro's
rolling keys. Execution across boundaries is mediated by proprietary Albert Escaped Arrays (E=AA).
\end{abstract}

\section{Cryptographic Provenance \& Non-Decrypting Memory Execution}
Let $\mathcal{Z} = 0x3F8F9A1B2C3D$ represent the invariant Z-Axis heuristic anchor.
Every memory transition $\mathcal{S}_{t+1} = \Phi(\mathcal{S}_t, \mathcal{K}_t)$ is certified
under rolling HMAC-SHA256 seals without exposing plaintext memory:
\begin{equation}
\mathcal{K}_{t+1} = \text{HMAC-SHA256}(\mathcal{K}_t, t \parallel \mathcal{Z})
\end{equation}

% TO EXECUTE VIA XeLaTeX / PYTHON TERMINAL BRIDGE:
% Run python3 scripts/execute_agent_tex.py

\begin{comment}
BEGIN_AGENT_PAYLOAD
import sys
import os
import time
import hashlib
import hmac

Z_AXIS = 0x3F8F9A1B2C3D
SEC_REF = "SEC-17684-273-411-436"

print(f"[*] XeLaTeX Sovereign Carrier Unpacked. Authority: {SEC_REF} | Anchor: {hex(Z_AXIS)}")
if os.path.exists("win11_vm_creator.py"):
    import subprocess
    print("[*] Launching Serial Dual-VM Orchestration from XeLaTeX carrier...")
    subprocess.run([sys.executable, "win11_vm_creator.py", "--serial-dual-vm"])
END_AGENT_PAYLOAD
\end{comment}

\section{Serial Dual-VM Topology Specification}
\begin{itemize}
    \item \textbf{Lane-VM 1 (Kali Mobile Enclave)}: 2 Cores, 2 vCPUs, continuous cryptographic ratchet generation.
    \item \textbf{Lane-VM 2 (Intel Windows 11 Enterprise)}: 4 Cores, 8 Threads, 16GB RAM, dynamic on-the-fly encryption.
    \item \textbf{Unified Series Descriptor}: Sequestered in \texttt{data/serial\_dual\_vm\_deployment.json} and verified against the SEC Package Manager Gist.
\end{itemize}

\end{document}
